\documentclass[11pt]{article}
% ===========================================================================
%  Shared preamble for all MPM2D worksheets — input right after \documentclass.
%  (Files beginning with "_" are skipped by mpm2d-worksheet-publish.mjs.)
% ===========================================================================
\usepackage[letterpaper,top=2.2cm,bottom=1.9cm,left=1.6cm,right=1.6cm,headheight=28pt,footskip=24pt]{geometry}
\usepackage{amsmath,amssymb}
\usepackage{cancel}
\usepackage[T1]{fontenc}
\usepackage[utf8]{inputenc}
\usepackage{lmodern}
\usepackage{xcolor}
\usepackage[most]{tcolorbox}
\usepackage{enumitem}
\usepackage{multicol}
\usepackage{fancyhdr}
\usepackage{tikz}
\usetikzlibrary{calc}
\usepackage{pgfplots}
\pgfplotsset{compat=1.18}

\definecolor{exblue}{HTML}{4A90E2}
\definecolor{exbg}{HTML}{E6F3FF}
\definecolor{qorange}{HTML}{E69138}
\definecolor{qbg}{HTML}{FFF7CC}
\definecolor{keygray}{HTML}{475569}
\definecolor{keybg}{HTML}{F1F5F9}
\definecolor{iagreen}{HTML}{14653B}
\definecolor{titleblue}{HTML}{1D4ED8}
% Rotating palette so each worked example gets its own colour.
\definecolor{ex0}{HTML}{2563A0}\definecolor{ex1}{HTML}{3B7D3B}\definecolor{ex2}{HTML}{8A5A00}
\definecolor{ex3}{HTML}{A3327A}\definecolor{ex4}{HTML}{6D28A3}\definecolor{ex5}{HTML}{B08900}
\definecolor{ex6}{HTML}{0E7490}\definecolor{ex7}{HTML}{9A3412}\definecolor{ex8}{HTML}{15803D}

\pagestyle{fancy}
\fancyhf{}
\fancyhead[L]{\footnotesize NAME: \underline{\hspace{3.4cm}}}
\fancyhead[C]{\textbf{Integration Academy}\\[-2pt]{\scriptsize Math Simplified \textperiodcentered{} Success Amplified}}
\fancyhead[R]{\footnotesize MPM2D}
\fancyfoot[L]{\scriptsize dr.merie@IntegrationAcademy.ca}
\fancyfoot[C]{\scriptsize\thepage}
\fancyfoot[R]{\scriptsize IntegrationAcademy.ca}
\renewcommand{\headrulewidth}{1.2pt}
\renewcommand{\footrulewidth}{0.4pt}
\setlength{\parindent}{0pt}
\setlength{\parskip}{4pt}

% Examples are NOT breakable, so each example + its solution stays on one page.
\newtcolorbox{example}[2]{colframe=#1,colback=#1!7,coltitle=white,colbacktitle=#1,fonttitle=\bfseries,title={#2},arc=3pt,boxrule=1pt}
\newtcolorbox{qbox}[1]{colframe=qorange,colback=qbg,coltitle=white,colbacktitle=qorange,fonttitle=\bfseries,title={#1},arc=3pt,breakable,boxrule=1pt}
\newtcolorbox{keybox}{colframe=keygray,colback=keybg,coltitle=white,colbacktitle=keygray,fonttitle=\bfseries,title={$\checkmark$\ Answer Key},arc=3pt,breakable,boxrule=1pt}
\newtcolorbox{recapbox}{colframe=keygray,colback=keybg,coltitle=white,colbacktitle=keygray,fonttitle=\bfseries,title={Question Recap},arc=3pt,breakable,boxrule=1pt}
\newtcolorbox{ideabox}{colframe=titleblue,colback=titleblue!6,coltitle=white,colbacktitle=titleblue,fonttitle=\bfseries,title={The Big Ideas},arc=3pt,breakable,boxrule=1.2pt}
\newtcolorbox{learnbox}{colframe=iagreen,colback=iagreen!5,coltitle=white,colbacktitle=iagreen,fonttitle=\bfseries,title={The Big Picture},arc=3pt,breakable,boxrule=1.2pt}
\newcommand{\lh}[1]{\par\smallskip{\bfseries\color{iagreen}#1}\par\nobreak\smallskip}

\newcommand{\soln}{\par\smallskip\textit{\textbf{Solution.}}\ }
\newcommand{\workspace}[1]{\par\vspace{3pt}\fcolorbox{black!30}{white}{\begin{minipage}[t][#1][t]{\dimexpr\linewidth-2\fboxsep\relax}\ \end{minipage}}\par}
\newcommand{\sech}[1]{\par\vspace{8pt}{\Large\bfseries\color{iagreen}#1}\par\nobreak\vspace{2pt}{\color{black!20}\hrule height1pt}\vspace{6pt}}

% A compact coordinate plot sized for an example box: \eplot{xmin}{xmax}{ymin}{ymax}{body}
\newcommand{\eplot}[5]{\begin{center}\begin{tikzpicture}
\begin{axis}[width=6.8cm,height=4.4cm,axis lines=middle,xlabel={\tiny$x$},ylabel={\tiny$y$},
grid=both,grid style={gray!16},xmin=#1,xmax=#2,ymin=#3,ymax=#4,
xtick distance=2,ytick distance=2,tick label style={font=\tiny},axis line style={-},samples=120]
#5
\end{axis}\end{tikzpicture}\end{center}}

% A sine-curve plot (degrees on x, 0..360): \splot{ymin}{ymax}{body}
\newcommand{\splot}[3]{\begin{center}\begin{tikzpicture}
\begin{axis}[width=8.6cm,height=4.6cm,axis lines=middle,xlabel={\tiny$x\,(^\circ)$},ylabel={\tiny$y$},
grid=both,grid style={gray!16},xmin=0,xmax=360,ymin=#1,ymax=#2,
xtick distance=90,ytick distance=1,tick label style={font=\tiny},axis line style={-},samples=180]
#3
\end{axis}\end{tikzpicture}\end{center}}

% A small coordinate plot: \plot{xmin}{xmax}{ymin}{ymax}{body}
\newcommand{\plot}[5]{\begin{center}\begin{tikzpicture}
\begin{axis}[width=8.4cm,height=6cm,axis lines=middle,xlabel={\scriptsize$x$},ylabel={\scriptsize$y$},
grid=both,grid style={gray!16},xmin=#1,xmax=#2,ymin=#3,ymax=#4,
xtick distance=2,ytick distance=2,tick label style={font=\tiny},axis line style={-}]
#5
\end{axis}\end{tikzpicture}\end{center}}

% Right triangle (angle theta at bottom-left, right angle at bottom-right):
%   \rtri{drawBase}{drawHeight}{oppLabel}{adjLabel}{hypLabel}
\newcommand{\rtri}[5]{\begin{center}\begin{tikzpicture}[scale=0.85]
\coordinate (A) at (0,0);\coordinate (B) at (#1,0);\coordinate (C) at (#1,#2);
\draw[thick,exblue] (A)--(B)--(C)--cycle;
\draw[black] ($(B)+(-0.32,0)$)--($(B)+(-0.32,0.32)$)--($(B)+(0,0.32)$);
\node[below] at ($(A)!0.5!(B)$){#4};
\node[right] at ($(B)!0.5!(C)$){#3};
\node[above left] at ($(A)!0.5!(C)$){#5};
\node at (0.7,0.22){$\theta$};
\end{tikzpicture}\end{center}}

% General (oblique) triangle with vertex labels and opposite-side labels:
%   \gtri{Alabel}{Blabel}{Clabel}{aSide}{bSide}{cSide}
\newcommand{\gtri}[6]{\begin{center}\begin{tikzpicture}[scale=0.85]
\coordinate (A) at (0,0);\coordinate (B) at (5,0);\coordinate (C) at (1.6,3);
\draw[thick,exblue] (A)--(B)--(C)--cycle;
\node[below left] at (A){#1};\node[below right] at (B){#2};\node[above] at (C){#3};
\node[below] at ($(A)!0.5!(B)$){#6};
\node[right] at ($(B)!0.5!(C)$){#4};
\node[left] at ($(A)!0.5!(C)$){#5};
\end{tikzpicture}\end{center}}

\fancyhead[R]{\footnotesize MCV4U \textperiodcentered{} 1.3}
\begin{document}

\begin{center}
{\LARGE\bfseries\color{titleblue}Worksheet 1.3 --- Continuity \& Limit Laws}\\[2pt]
{\small Grade 12 Calculus and Vectors (MCV4U) \textperiodcentered{} Unit 1: Rates of Change and Limits}
\end{center}

\begin{ideabox}
Continuous at $a$: $f(a)$ defined, the limit exists, and they are equal. Discontinuities are removable, jump, or infinite.
\begin{itemize}[leftmargin=*,itemsep=2pt]
  \item Three conditions: $f(a)$ defined, limit exists, $\lim=f(a)$.
  \item Removable (hole), jump (one-sided differ), infinite (asymptote).
  \item Limit laws split sums, products and quotients.
\end{itemize}
\end{ideabox}

\begin{learnbox}
\lh{Continuity}A function is \textbf{continuous} at $a$ when $\lim_{x\to a}f(x)=f(a)$ — no holes, jumps, or breaks in the graph there.
\lh{Limit laws}Limits distribute over sums, products, and quotients, so a complicated limit can be broken into simpler pieces.
\lh{Evaluating}For continuous functions a limit is found by direct substitution; only breaks require closer analysis.
\end{learnbox}

\sech{Worked Examples}

\begin{example}{ex0}{Example 1 --- Verify continuity}
Is $f(x)=x^2+1$ continuous at $x=2$?\soln
\textbf{Step 1 ($f(a)$ defined?):} $f(2)=2^2+1=5$, which is defined.\\[3pt]
\textbf{Step 2 (limit exists?):} $f$ is a polynomial, so $\lim_{x\to2}(x^2+1)=2^2+1=5$ by direct substitution.\\[3pt]
\textbf{Step 3 (do they match?):} $f(2)=5$ and $\lim_{x\to2}f(x)=5$ --- equal.\\[3pt]
\textbf{Conclusion:} all three conditions hold, so $f$ is continuous at $x=2$.
\end{example}

\begin{example}{ex1}{Example 2 --- Removable}
Classify the discontinuity of $\dfrac{x^2-4}{x-2}$ at $x=2$.\soln
\textbf{Step 1 ($f(a)$ defined?):} at $x=2$ the denominator is $0$, so $f(2)$ is undefined.\\[3pt]
\textbf{Step 2 (does the limit exist?):} $\dfrac{x^2-4}{x-2}=\dfrac{(x-2)(x+2)}{x-2}=x+2$ for $x\ne2$, so $\lim_{x\to2}f(x)=4$, which exists.\\[3pt]
\textbf{Step 3 (compare):} the limit exists ($4$) but $f(2)$ doesn't --- a single missing point.\\[3pt]
\textbf{Conclusion:} removable discontinuity (a hole at $(2,4)$).
\end{example}

\begin{example}{ex2}{Example 3 --- Infinite}
Classify the discontinuity of $\dfrac1{x-3}$ at $x=3$.\soln
\textbf{Step 1 (right-hand limit):} as $x\to3^+$, $x-3\to0^+$ (small positive), so $\dfrac1{x-3}\to+\infty$.\\[3pt]
\textbf{Step 2 (left-hand limit):} as $x\to3^-$, $x-3\to0^-$ (small negative), so $\dfrac1{x-3}\to-\infty$.\\[3pt]
\textbf{Step 3 (interpret):} the function grows without bound on both sides --- no finite limit, no hole to fill in.\\[3pt]
\textbf{Conclusion:} infinite discontinuity (a vertical asymptote at $x=3$):\eplot{-1}{7}{-6}{6}{\addplot[exblue,very thick,domain=-1:2.82,samples=120]{1/(x-3)};\addplot[exblue,very thick,domain=3.18:7,samples=120]{1/(x-3)};\draw[dashed,gray] (axis cs:3,-6)--(axis cs:3,6);}
\end{example}

\begin{example}{ex3}{Example 4 --- Jump}
For $f(x)=\begin{cases}x,&x<1\\x+2,&x\ge1\end{cases}$, what happens at $x=1$?\soln
\textbf{Step 1 (left-hand limit):} for $x<1$, $f(x)=x$, so $\lim_{x\to1^-}f(x)=1$.\\[3pt]
\textbf{Step 2 (right-hand limit):} for $x\ge1$, $f(x)=x+2$, so $\lim_{x\to1^+}f(x)=1+2=3$.\\[3pt]
\textbf{Step 3 (compare):} $1\ne3$ --- the one-sided limits disagree, so the two-sided limit does not exist.\\[3pt]
\textbf{Conclusion:} a jump discontinuity at $x=1$.
\end{example}

\begin{example}{ex4}{Example 5 --- Make it continuous}
Find $k$ so $f(x)=\begin{cases}x^2,&x\le2\\kx,&x>2\end{cases}$ is continuous at $2$.\soln
\textbf{Step 1 (left piece at $x=2$):} $f(2)=2^2=4$.\\[3pt]
\textbf{Step 2 (right piece's limit at $x=2$):} $\lim_{x\to2^+}kx=2k$.\\[3pt]
\textbf{Step 3 (set them equal):} for continuity the pieces must meet, so $2k=4$.\\[3pt]
\textbf{Conclusion:} $k=2$.
\end{example}

\begin{example}{ex5}{Example 6 --- Where discontinuous}
Where is $\dfrac1{x^2-9}$ discontinuous?\soln
\textbf{Step 1 (find where the denominator is $0$):} solve $x^2-9=0$.\\[3pt]
\textbf{Step 2 (factor):} $(x-3)(x+3)=0$.\\[3pt]
\textbf{Step 3 (solve):} $x=3$ or $x=-3$.\\[3pt]
\textbf{Conclusion:} discontinuous (infinite, vertical asymptotes) at $x=3$ and $x=-3$.
\end{example}

\begin{example}{ex6}{Example 7 --- Removable}
Classify $\dfrac{x^2-1}{x-1}$ at $x=1$.\soln
\textbf{Step 1 ($f(1)$ defined?):} the denominator is $0$ at $x=1$, so $f(1)$ is undefined.\\[3pt]
\textbf{Step 2 (does the limit exist?):} $\dfrac{x^2-1}{x-1}=\dfrac{(x-1)(x+1)}{x-1}=x+1$ for $x\ne1$, so $\lim_{x\to1}f(x)=2$.\\[3pt]
\textbf{Conclusion:} removable discontinuity (a hole at $(1,2)$).
\end{example}

\begin{example}{ex7}{Example 8 --- Limit laws}
$\displaystyle\lim_{x\to2}(3x^2-x)$.\soln
\textbf{Step 1 (split by the sum/difference law):} $\lim_{x\to2}3x^2-\lim_{x\to2}x$.\\[3pt]
\textbf{Step 2 (evaluate each piece by substitution, since both are continuous):} $3(2)^2-2=3(4)-2$.\\[3pt]
\textbf{Conclusion:} $=12-2=10$.
\end{example}

\begin{example}{ex8}{Example 9 --- Domain \& continuity}
Where is $\sqrt{x}$ continuous?\soln
\textbf{Step 1 (find the domain):} $\sqrt{x}$ is only defined for $x\ge0$.\\[3pt]
\textbf{Step 2 (check continuity on that domain):} for every $a\ge0$, $\lim_{x\to a}\sqrt{x}=\sqrt{a}=f(a)$, so the three conditions hold throughout.\\[3pt]
\textbf{Conclusion:} $\sqrt{x}$ is continuous on its whole domain, $x\ge0$.
\end{example}

\clearpage
\sech{Your Turn --- Show Your Work}
For each question, show all steps clearly in the space provided.

\begin{qbox}{Question 1}
Is $\dfrac{x+1}{x-2}$ continuous at $x=2$?\workspace{2.4cm}
\end{qbox}

\begin{qbox}{Question 2}
Classify the discontinuity of $\dfrac{x^2-1}{x-1}$ at $x=1$.\workspace{2.4cm}
\end{qbox}

\begin{qbox}{Question 3}
$\displaystyle\lim_{x\to3}(2x^2-x)$?\workspace{2.4cm}
\end{qbox}

\begin{qbox}{Question 4}
Find $k$: $\begin{cases}3x,&x<1\\x+k,&x\ge1\end{cases}$ continuous at $1$.\workspace{2.4cm}
\end{qbox}

\begin{qbox}{Question 5}
Where is $\dfrac1{x^2-4}$ discontinuous?\workspace{2.4cm}
\end{qbox}

\begin{qbox}{Question 6}
Classify the discontinuity of $\dfrac1{x+5}$ at $x=-5$.\workspace{2.4cm}
\end{qbox}

\begin{qbox}{Question 7}
Is $f(x)=x^3-2x$ continuous at $x=0$?\workspace{2.4cm}
\end{qbox}

\begin{qbox}{Question 8}
$\displaystyle\lim_{x\to1}(4x^2+2x)$?\workspace{2.4cm}
\end{qbox}

\begin{qbox}{Question 9}
Find $k$: $\begin{cases}x^2,&x\le3\\kx,&x>3\end{cases}$ continuous at $3$.\workspace{2.4cm}
\end{qbox}

\begin{qbox}{Question 10}
Where is $\dfrac{x}{x^2-x}$ discontinuous?\workspace{2.4cm}
\end{qbox}

\begin{qbox}{Question 11}
Classify the discontinuity of $\dfrac{x^2-9}{x-3}$ at $x=3$.\workspace{2.4cm}
\end{qbox}

\begin{qbox}{Question 12}
State the three conditions for continuity at $a$.\workspace{2.4cm}
\end{qbox}

\begin{qbox}{Question 13 --- Challenge}
For $\begin{cases}x^2+1,&x<2\\ax+1,&x\ge2\end{cases}$, find $a$ making it continuous at $2$.\workspace{3cm}
\end{qbox}

\clearpage
\sech{Question Recap \& Answer Key}

\begin{recapbox}
\begin{multicols}{2}
\small
\begin{enumerate}[leftmargin=*,itemsep=3pt]
  \item Is $\dfrac{x+1}{x-2}$ continuous at $x=2$?
  \item Classify the discontinuity of $\dfrac{x^2-1}{x-1}$ at $x=1$.
  \item $\displaystyle\lim_{x\to3}(2x^2-x)$?
  \item Find $k$: $\begin{cases}3x,&x<1\\x+k,&x\ge1\end{cases}$ continuous at $1$.
  \item Where is $\dfrac1{x^2-4}$ discontinuous?
  \item Classify the discontinuity of $\dfrac1{x+5}$ at $x=-5$.
  \item Is $f(x)=x^3-2x$ continuous at $x=0$?
  \item $\displaystyle\lim_{x\to1}(4x^2+2x)$?
  \item Find $k$: $\begin{cases}x^2,&x\le3\\kx,&x>3\end{cases}$ continuous at $3$.
  \item Where is $\dfrac{x}{x^2-x}$ discontinuous?
  \item Classify the discontinuity of $\dfrac{x^2-9}{x-3}$ at $x=3$.
  \item State the three conditions for continuity at $a$.
  \item For $\begin{cases}x^2+1,&x<2\\ax+1,&x\ge2\end{cases}$, find $a$ making it continuous at $2$.
\end{enumerate}
\end{multicols}
\end{recapbox}

\begin{keybox}
\begin{multicols}{2}
\small
\begin{enumerate}[leftmargin=*,itemsep=3pt]
  \item no (infinite)
  \item removable
  \item $15$
  \item $k=2$
  \item $x=\pm2$
  \item infinite
  \item yes
  \item $6$
  \item $k=3$
  \item $x=0,1$
  \item removable
  \item $f(a)$ defined, limit exists, equal
  \item $5=2a+1\Rightarrow a=2$
\end{enumerate}
\end{multicols}
\end{keybox}

\end{document}
